To confirm and apply our findings to real-world problems, we simulate the heat treament of a Ti-Al alloy material after the direct energy deposition.

\subsection{Model setup and simulation details}
The simulation cell is defined as in \cref{fig:heat_treatment}. In this simulation cell, the top \SI{1}{\nano\meter}-thick layer of Ti and the bottom \SI{1}{\nano\meter}-thick layer of Al are fixed such that they are not allowed to move freely in the xy-plane while the movement along the z-axis is not constrained. A \SI{1}{\nano\meter}-wide margin of the TiAl\textsubscript{3} is fixed in all three direction, preventing the TiAl\textsubscript{3} layer (and thus the entire system) from floating freely in the space.

\begin{figure}[htbp]
    \centering
    \includegraphics[width=0.6\textwidth]{figs/heat_treatment_sys}
    \caption{The simulation cell for the heat treatment process. Besides the dimensions shown in the figure, the thicknesses of the Al, TiAl\textsubscript{3} and Ti layers are \SI{10}{\nano\meter}, \SI{4}{\nano\meter} and \SI{4}{\nano\meter} respectively. The TiAl\textsubscript{3} layer consists of $4$ randomly oriented grains, forming grain boundaries. The lattice orientations of the Al and Ti layers are also random. (The supporting information includes all lattice oritentations.)}
    \label{fig:heat_treatment}
\end{figure}

The system is first equilibrated at \SI{300}{\kelvin} for \SI{1}{\nano\meter}. Then, it is heated up to \SI{895.15}{\kelvin} in \SI{300}{\pico\second}, corresponding to a heating rate of approximately \SI{2}{\kelvin/\pico\second}. Next, the system is kept at \SI{895.15}{\kelvin} for \SI{30}{\nano\second}. \SI{895.15}{\kelvin} falls in the temperature range identified in several experimental studies for Ti-Al intermetallic system annealing\cite{bizot2020reactivity, mirjalili2013kinetics, xu2006growth} and is a typical temperature used in this study\cite{mirjalili2013kinetics}. In the system of interest, at this temperature, the dominating dynamics is the diffusion of Al atoms from the pure Al side through the TiAl\textsubscript{3} grain boundaries to the Ti side. So, we analyze the evolution of the system by focusing on the diffusion of the Al atoms that originates from the pure Al layer (as opposed to the Al atoms originally in the TiAl\textsubscript{3}).